% !TeX program = tectonic
\documentclass[11pt,a4paper]{article}
\usepackage{fontspec}
\usepackage[english]{babel}
\babelfont{rm}[Language=Default]{Liberation Serif}
\babelfont[Language=Default]{sf}{Carlito}
\babelfont[Language=Default]{tt}{DejaVu Sans Mono}
\usepackage{amsmath,amssymb,amsthm}
\usepackage{geometry}
\geometry{a4paper, top=2.4cm, bottom=2.4cm, left=2.4cm, right=2.4cm}
\usepackage{booktabs,tabularx,enumitem,xcolor,tcolorbox}
\tcbuselibrary{breakable,skins}
\definecolor{linkblue}{HTML}{1F4E79}
\definecolor{statgreen}{HTML}{2F6B3A}
\definecolor{goldacc}{HTML}{8B7E5A}
\usepackage[colorlinks=true,linkcolor=linkblue,urlcolor=linkblue,unicode=true]{hyperref}
\theoremstyle{plain}
\newtheorem{theorem}{Theorem}
\newtheorem{lemma}[theorem]{Lemma}
\newtheorem{proposition}[theorem]{Proposition}
\newtheorem{corollary}[theorem]{Corollary}
\theoremstyle{definition}
\newtheorem{definition}[theorem]{Definition}
\theoremstyle{remark}
\newtheorem{remark}[theorem]{Remark}
\newtcolorbox{statusbox}[1][]{colback=statgreen!5,colframe=statgreen!75,
  title={\textbf{Summary of results:} #1},fonttitle=\small,boxrule=0.7pt,arc=2pt,breakable}
\newtcolorbox{databox}[1][]{colback=goldacc!6,colframe=goldacc!80,
  title={\textbf{Protocol data:} #1},fonttitle=\small,boxrule=0.7pt,arc=2pt,breakable}
\newtcolorbox{verifybox}[1][]{colback=linkblue!5,colframe=linkblue!80,
  title={\textbf{Verification:} #1},fonttitle=\small,boxrule=0.7pt,arc=2pt,breakable}
\widowpenalty=10000
\clubpenalty=10000
\setlength{\parskip}{0.35em}
\newcommand{\CC}{\mathbb{C}}
\newcommand{\RR}{\mathbb{R}}
\newcommand{\ZZ}{\mathbb{Z}}
\newcommand{\QQ}{\mathbb{Q}}
\newcommand{\Ga}{\Gamma}
\newcommand{\Bfun}{\mathrm{B}}
\newcommand{\im}{\operatorname{Im}}
\newcommand{\re}{\operatorname{Re}}
\newcommand{\lcm}{\operatorname{lcm}}
\newcommand{\SNF}{\operatorname{SNF}}
\newcommand{\Jac}{\operatorname{Jac}}
\newcommand{\rank}{\operatorname{rank}}
\newcommand{\Ind}{\operatorname{Index}}
\newcommand{\disc}{\operatorname{disc}}
\begin{document}
\begin{center}
{\LARGE\bfseries Certificate G: Hodge Lattice Indices}\\[0.4em]
{\large lattice spectra, CM fractions, and the Chowla–Selberg layer}\\[0.6em]
{\normalsize Isaev Iskhak Khamzatovich}\\[0.2em]
{\normalsize The <<Dynamic Principle>> Program --- the \texttt{hodge-laboratory}}\\[0.15em]
{\small Standalone theorem monograph · Монография 15/16}
\end{center}
\vspace{0.4em}
{\color{linkblue}\hrule height 0.8pt}
\vspace{0.8em}

\section{Setting}

Certificate G provides the universal normalization of the spectrum
of an arbitrary lattice and computes Hodge lattice indices in
number fields --- the bridge from integer geometry to the
$\Ga$-levels of certificate H.

\begin{theorem}[G1--G6]\label{th:Gstand}
\begin{enumerate}[leftmargin=2em]
\item[(G1)] Universal spectrum:
$\lambda_{m,n}=(2\pi)^2|m\tau-n|^2/(\im\tau)^2$; the factor
$4\pi^2$ is the universal Tate twist.
\item[(G2)] For reduced lattices $\lambda_1=4\pi^2/(\im\tau)^2$;
for CM fields --- exact fractions: the families $16\pi^2/d$ and
$4\pi^2/d$. The levels $\pi/15$ and $\pi/30$ are spectral units of
$\ZZ[\sqrt{-15}]$, $\ZZ[\sqrt{-30}]$: $\lambda_1/(4\pi)=\pi/N$.
\item[(G3)] Indices in number fields are recovered with uniqueness
$2\varepsilon<1/Q$ (field-LLL with the triple $(p,q,s)$).
\item[(G4)] Chowla--Selberg layer: $\omega(i)^2=\Ga(1/4)^4/(16\pi)$;
$\omega(\tau_7)^2=\frac{21+\sqrt{-7}}{896}
\bigl[\Ga(\frac17)\Ga(\frac27)\Ga(\frac47)\bigr]^2/\pi^2$;
$\omega(\tau_3)^2=\frac{3-\sqrt{-3}}8
\Ga(\frac13)^6/(\pi^2 2^{8/3})$; $d=15$: the $j$-pair in
$\QQ(\sqrt5)$ and $R=\frac{(2\sqrt5-3)+(\sqrt5-2)\sqrt{-3}}2$.
\item[(G5)] K3 bridge: the quadrant integral
$\Ga(1/4)^4/(16\pi\sqrt2)$; the index $\sqrt2/32$ in the
$\Ga$-scale.
\item[(G6)] The $\mu_N$-quantum ladder: quantum $2\pi/N$, spin half
$\pi/N$; $\delta_C=\pi/7$ is the $N=7$ rung; $\pi/15$, $\pi/30$ ---
the next rungs.
\end{enumerate}
\end{theorem}

\begin{proof}
(G1) The flat Laplacian on the elliptic curve with lattice
$\ZZ\tau+\ZZ$: eigenvalues $\propto|m\tau-n|^2$; normalization by
$(2\pi)^2$.
(G2) For a reduced $\tau$ the minimum $|m\tau-n|^2=1$ is attained at
the shortest vector; the CM case: $\tau=\sqrt{-d}$ or
$\frac{1+\sqrt{-d}}2$, the fraction unfolds into the exact
$4\pi^2/d$ or $16\pi^2/d$ by the shortest-vector classification;
for $d=15$ the norm computation gives $\lambda_1/(4\pi)=\pi/15$.
(G3) The field separation is bounded below via the discriminant;
LLL in the basis $(p,q,s)$ is unique at $2\varepsilon<1/Q$; the
recovery residual $2.31\cdot10^{-102}$.
(G4) Chowla--Selberg: the squared period is a $\Ga$-monomial with
rational exponents; all identities verified with residuals
$<10^{-100}$.
(G5) The quadrant split and a change of variables give the exact
value; residual $3.92\cdot10^{-43}$.
(G6) The spin half follows from the double cover; the ladder agrees
with the stand triples.
\end{proof}

\begin{databox}[numbers]
The non-CM test: the minimal polynomial $4X^3-1$ --- field-LLL
distinguishes CM from non-CM. Field-LLL $d=15$: the recovery of
$\frac{3+5\sqrt{-15}}8$. The spectral fractions confirmed on the
nine G1 lattices; the families $16\pi^2/d$, $4\pi^2/d$ --- exact.
\end{databox}

\begin{statusbox}[summary]
Certificate G accepted ($6/6$): the universal normalization, exact
fractionality, field-LLL, the Chowla--Selberg layer, the K3 bridge,
the ladder --- all points certified.
\end{statusbox}

\begin{verifybox}[reproduction]
\texttt{python3 laboratory.py} $\to$ ``Certificates $\to$ G''; the
designer: an arbitrary $d$ --- the lab computes $\lambda_1$ and
prints the fraction.
\end{verifybox}

\vfill
{\color{linkblue}\hrule height 0.6pt}
\vspace{0.5em}
{\small\color{gray}
License: individual exclusive license of Isaev Iskhak Khamzatovich.
All rights to the computations, formulas, and texts belong to the author.
Verification: \texttt{python3 laboratory.py} (``Stands''/``Certificates''),
Lean: \texttt{verification/lean/HodgeLaboratory.lean}.
}
\end{document}
